\documentclass[10pt,letterpaper,sans]{moderncv}

\moderncvstyle{banking}
\moderncvcolor{blue}

\usepackage[scale=0.88]{geometry}
\setlength{\hintscolumnwidth}{2.1cm}
\renewcommand*{\namefont}{\fontsize{24}{26}\bfseries\upshape}
% \renewcommand*{\titlefont}{\Large\mdseries\upshape}
% \renewcommand*{\addressfont}{\small\mdseries\upshape}

\name{Max}{Siegel}
\title{Human-Centered AI}
\address{Cambridge, MA}{}{}
\email{maxs@mit.edu}

\begin{document}
\makecvtitle
\vspace{-1.1em}

\section{Profile}
\cvitem{}{Human-centered AI researcher with a background in computational cognitive science, perception, and social reasoning. Builds computational models of interpretation, inference, and learning; designs human-subject studies to evaluate those models; and studies how cognition can inform intelligent systems and interaction design.}

\section{Experience}
\cventry{2018--Present}{Research Scientist}{MIT, Department of Brain and Cognitive Sciences}{Cambridge, MA}{}{
\begin{itemize}
\item Conduct interdisciplinary research on how people perceive, search, infer, communicate, and learn under uncertainty in MIT's Computational Cognitive Science group.
\item Develop computational models of perception, reasoning, and interaction, including analysis-by-synthesis and theory-driven accounts of human inference.
\item Design behavioral studies and quantitative evaluations that connect human performance to computational models and intelligent-system behavior.
\item Collaborate across cognitive science, computer science, neuroscience, and developmental psychology on projects involving visual, auditory, linguistic, and social signals.
\item Publish research on cross-cultural communication, common ground, social interaction, multimodal inference, and exploratory learning.
\item Built a publication record of 23 papers, including 4 co-first-authored papers and articles in \textit{Nature Human Behaviour}, \textit{Nature Communications}, and other interdisciplinary venues.
\end{itemize}}

\section{Strengths}
\cvitem{Modeling}{Computational modeling, analysis-by-synthesis, hypothesis testing, quantitative evaluation}
\cvitem{Domains}{Human-AI interaction, social reasoning, common ground, communication, multimodal perception, concept learning}
\cvitem{Human Research}{Behavioral experiments, psychophysics, study design, human-subject evaluation, interdisciplinary collaboration}

\section{Selected Impact}
\cvitem{}{Co-first-authored \textit{Nature Human Behaviour} research on physically grounded 3D perception using analysis-by-synthesis and neural perceptual models.}
\cvitem{}{Co-first-authored \textit{Nature Communications} research on exploratory learning as adaptive, precisely calibrated hypothesis testing.}
\cvitem{}{Published many computational and behavioral studies of human perception, concept learning,  social interaction, and multimodal reasoning from visual and auditory evidence.}

\section{Selected Publications}
\cvitem{2023}{Yildirim, I., \textbf{Siegel, M. H.}, Soltani, A. A., Ray Chaudhuri, S., \& Tenenbaum, J. B. Perception of 3D shape integrates intuitive physics and analysis-by-synthesis. \textit{Nature Human Behaviour}.}
\cvitem{2021}{\textbf{Siegel, M. H.}, Magid, R. W., Pelz, M., Tenenbaum, J. B., \& Schulz, L. E. Children's exploratory play tracks the discriminability of hypotheses. \textit{Nature Communications}.}

\section{Education}
\cventry{2018}{Ph.D. in Brain and Cognitive Sciences}{MIT}{Cambridge, MA}{}{Graduated September 2018.}

\end{document}
